% Options for packages loaded elsewhere
% Options for packages loaded elsewhere
\PassOptionsToPackage{unicode}{hyperref}
\PassOptionsToPackage{hyphens}{url}
%
\documentclass[
  ignorenonframetext,
  aspectratio=169,
]{beamer}
\newif\ifbibliography
\usepackage{pgfpages}
\setbeamertemplate{caption}[numbered]
\setbeamertemplate{caption label separator}{: }
\setbeamercolor{caption name}{fg=normal text.fg}
\beamertemplatenavigationsymbolsempty
% remove section numbering
\setbeamertemplate{part page}{
  \centering
  \begin{beamercolorbox}[sep=16pt,center]{part title}
    \usebeamerfont{part title}\insertpart\par
  \end{beamercolorbox}
}
\setbeamertemplate{section page}{
  \centering
  \begin{beamercolorbox}[sep=12pt,center]{section title}
    \usebeamerfont{section title}\insertsection\par
  \end{beamercolorbox}
}
\setbeamertemplate{subsection page}{
  \centering
  \begin{beamercolorbox}[sep=8pt,center]{subsection title}
    \usebeamerfont{subsection title}\insertsubsection\par
  \end{beamercolorbox}
}
% Prevent slide breaks in the middle of a paragraph
\widowpenalties 1 10000
\raggedbottom
\AtBeginPart{
  \frame{\partpage}
}
\AtBeginSection{
  \ifbibliography
  \else
    \frame{\sectionpage}
  \fi
}
\AtBeginSubsection{
  \frame{\subsectionpage}
}
\usepackage{iftex}
\ifPDFTeX
  \usepackage[T1]{fontenc}
  \usepackage[utf8]{inputenc}
  \usepackage{textcomp} % provide euro and other symbols
\else % if luatex or xetex
  \usepackage{unicode-math} % this also loads fontspec
  \defaultfontfeatures{Scale=MatchLowercase}
  \defaultfontfeatures[\rmfamily]{Ligatures=TeX,Scale=1}
\fi
\usepackage{lmodern}

\ifPDFTeX\else
  % xetex/luatex font selection
\fi
% Use upquote if available, for straight quotes in verbatim environments
\IfFileExists{upquote.sty}{\usepackage{upquote}}{}
\IfFileExists{microtype.sty}{% use microtype if available
  \usepackage[]{microtype}
  \UseMicrotypeSet[protrusion]{basicmath} % disable protrusion for tt fonts
}{}
\makeatletter
\@ifundefined{KOMAClassName}{% if non-KOMA class
  \IfFileExists{parskip.sty}{%
    \usepackage{parskip}
  }{% else
    \setlength{\parindent}{0pt}
    \setlength{\parskip}{6pt plus 2pt minus 1pt}}
}{% if KOMA class
  \KOMAoptions{parskip=half}}
\makeatother

\usepackage{color}
\usepackage{fancyvrb}
\newcommand{\VerbBar}{|}
\newcommand{\VERB}{\Verb[commandchars=\\\{\}]}
\DefineVerbatimEnvironment{Highlighting}{Verbatim}{commandchars=\\\{\}}
% Add ',fontsize=\small' for more characters per line
\usepackage{framed}
\definecolor{shadecolor}{RGB}{241,243,245}
\newenvironment{Shaded}{\begin{snugshade}}{\end{snugshade}}
\newcommand{\AlertTok}[1]{\textcolor[rgb]{0.68,0.00,0.00}{#1}}
\newcommand{\AnnotationTok}[1]{\textcolor[rgb]{0.37,0.37,0.37}{#1}}
\newcommand{\AttributeTok}[1]{\textcolor[rgb]{0.40,0.45,0.13}{#1}}
\newcommand{\BaseNTok}[1]{\textcolor[rgb]{0.68,0.00,0.00}{#1}}
\newcommand{\BuiltInTok}[1]{\textcolor[rgb]{0.00,0.23,0.31}{#1}}
\newcommand{\CharTok}[1]{\textcolor[rgb]{0.13,0.47,0.30}{#1}}
\newcommand{\CommentTok}[1]{\textcolor[rgb]{0.37,0.37,0.37}{#1}}
\newcommand{\CommentVarTok}[1]{\textcolor[rgb]{0.37,0.37,0.37}{\textit{#1}}}
\newcommand{\ConstantTok}[1]{\textcolor[rgb]{0.56,0.35,0.01}{#1}}
\newcommand{\ControlFlowTok}[1]{\textcolor[rgb]{0.00,0.23,0.31}{\textbf{#1}}}
\newcommand{\DataTypeTok}[1]{\textcolor[rgb]{0.68,0.00,0.00}{#1}}
\newcommand{\DecValTok}[1]{\textcolor[rgb]{0.68,0.00,0.00}{#1}}
\newcommand{\DocumentationTok}[1]{\textcolor[rgb]{0.37,0.37,0.37}{\textit{#1}}}
\newcommand{\ErrorTok}[1]{\textcolor[rgb]{0.68,0.00,0.00}{#1}}
\newcommand{\ExtensionTok}[1]{\textcolor[rgb]{0.00,0.23,0.31}{#1}}
\newcommand{\FloatTok}[1]{\textcolor[rgb]{0.68,0.00,0.00}{#1}}
\newcommand{\FunctionTok}[1]{\textcolor[rgb]{0.28,0.35,0.67}{#1}}
\newcommand{\ImportTok}[1]{\textcolor[rgb]{0.00,0.46,0.62}{#1}}
\newcommand{\InformationTok}[1]{\textcolor[rgb]{0.37,0.37,0.37}{#1}}
\newcommand{\KeywordTok}[1]{\textcolor[rgb]{0.00,0.23,0.31}{\textbf{#1}}}
\newcommand{\NormalTok}[1]{\textcolor[rgb]{0.00,0.23,0.31}{#1}}
\newcommand{\OperatorTok}[1]{\textcolor[rgb]{0.37,0.37,0.37}{#1}}
\newcommand{\OtherTok}[1]{\textcolor[rgb]{0.00,0.23,0.31}{#1}}
\newcommand{\PreprocessorTok}[1]{\textcolor[rgb]{0.68,0.00,0.00}{#1}}
\newcommand{\RegionMarkerTok}[1]{\textcolor[rgb]{0.00,0.23,0.31}{#1}}
\newcommand{\SpecialCharTok}[1]{\textcolor[rgb]{0.37,0.37,0.37}{#1}}
\newcommand{\SpecialStringTok}[1]{\textcolor[rgb]{0.13,0.47,0.30}{#1}}
\newcommand{\StringTok}[1]{\textcolor[rgb]{0.13,0.47,0.30}{#1}}
\newcommand{\VariableTok}[1]{\textcolor[rgb]{0.07,0.07,0.07}{#1}}
\newcommand{\VerbatimStringTok}[1]{\textcolor[rgb]{0.13,0.47,0.30}{#1}}
\newcommand{\WarningTok}[1]{\textcolor[rgb]{0.37,0.37,0.37}{\textit{#1}}}

\usepackage{longtable,booktabs,array}
\usepackage{calc} % for calculating minipage widths
\usepackage{caption}
% Make caption package work with longtable
\makeatletter
\def\fnum@table{\tablename~\thetable}
\makeatother
\usepackage{graphicx}
\makeatletter
\newsavebox\pandoc@box
\newcommand*\pandocbounded[1]{% scales image to fit in text height/width
  \sbox\pandoc@box{#1}%
  \Gscale@div\@tempa{\textheight}{\dimexpr\ht\pandoc@box+\dp\pandoc@box\relax}%
  \Gscale@div\@tempb{\linewidth}{\wd\pandoc@box}%
  \ifdim\@tempb\p@<\@tempa\p@\let\@tempa\@tempb\fi% select the smaller of both
  \ifdim\@tempa\p@<\p@\scalebox{\@tempa}{\usebox\pandoc@box}%
  \else\usebox{\pandoc@box}%
  \fi%
}
% Set default figure placement to htbp
\def\fps@figure{htbp}
\makeatother





\setlength{\emergencystretch}{3em} % prevent overfull lines

\providecommand{\tightlist}{%
  \setlength{\itemsep}{0pt}\setlength{\parskip}{0pt}}



 


\makeatletter
\@ifpackageloaded{tcolorbox}{}{\usepackage[skins,breakable]{tcolorbox}}
\@ifpackageloaded{fontawesome5}{}{\usepackage{fontawesome5}}
\definecolor{quarto-callout-color}{HTML}{909090}
\definecolor{quarto-callout-note-color}{HTML}{0758E5}
\definecolor{quarto-callout-important-color}{HTML}{CC1914}
\definecolor{quarto-callout-warning-color}{HTML}{EB9113}
\definecolor{quarto-callout-tip-color}{HTML}{00A047}
\definecolor{quarto-callout-caution-color}{HTML}{FC5300}
\definecolor{quarto-callout-color-frame}{HTML}{acacac}
\definecolor{quarto-callout-note-color-frame}{HTML}{4582ec}
\definecolor{quarto-callout-important-color-frame}{HTML}{d9534f}
\definecolor{quarto-callout-warning-color-frame}{HTML}{f0ad4e}
\definecolor{quarto-callout-tip-color-frame}{HTML}{02b875}
\definecolor{quarto-callout-caution-color-frame}{HTML}{fd7e14}
\makeatother
\makeatletter
\@ifpackageloaded{caption}{}{\usepackage{caption}}
\AtBeginDocument{%
\ifdefined\contentsname
  \renewcommand*\contentsname{Table of contents}
\else
  \newcommand\contentsname{Table of contents}
\fi
\ifdefined\listfigurename
  \renewcommand*\listfigurename{List of Figures}
\else
  \newcommand\listfigurename{List of Figures}
\fi
\ifdefined\listtablename
  \renewcommand*\listtablename{List of Tables}
\else
  \newcommand\listtablename{List of Tables}
\fi
\ifdefined\figurename
  \renewcommand*\figurename{Figure}
\else
  \newcommand\figurename{Figure}
\fi
\ifdefined\tablename
  \renewcommand*\tablename{Table}
\else
  \newcommand\tablename{Table}
\fi
}
\@ifpackageloaded{float}{}{\usepackage{float}}
\floatstyle{ruled}
\@ifundefined{c@chapter}{\newfloat{codelisting}{h}{lop}}{\newfloat{codelisting}{h}{lop}[chapter]}
\floatname{codelisting}{Listing}
\newcommand*\listoflistings{\listof{codelisting}{List of Listings}}
\makeatother
\makeatletter
\makeatother
\makeatletter
\@ifpackageloaded{caption}{}{\usepackage{caption}}
\@ifpackageloaded{subcaption}{}{\usepackage{subcaption}}
\makeatother

\usepackage{bookmark}
\IfFileExists{xurl.sty}{\usepackage{xurl}}{} % add URL line breaks if available
\urlstyle{same}
\hypersetup{
  pdftitle={Análise Fatorial},
  hidelinks,
  pdfcreator={LaTeX via pandoc}}



\title{Análise Fatorial}
\author{}
\date{}

\begin{document}
\frame{\titlepage}


\begin{frame}{Análise Fatorial}
\phantomsection\label{anuxe1lise-fatorial}
\begin{itemize}
\tightlist
\item
  \textbf{Objetivo:} Descrever as relações de covariância entre muitas
  variáveis em termos de poucas quantidades aleatórias subjacentes e não
  observáveis
\end{itemize}

\pause

\begin{itemize}
\tightlist
\item
  \textbf{Motivação:} Variáveis de um grupo altamente correlacionadas
  entre si, mas com pequenas correlações de outros grupos.

  \begin{itemize}
  \tightlist
  \item
    É concebível que cada grupo de variáveis represente um fator (ou
    construto) que seja o responsável pelas correlações observadas
  \end{itemize}
\end{itemize}
\end{frame}

\begin{frame}{Análise Fatorial}
\phantomsection\label{anuxe1lise-fatorial-1}
Vamos supor cinco variáveis \(Y_1\), \(Y_2\), \(Y_3\), \(Y_4\) e \(Y_5\)
correspondentes às notas de estudantes em Espanhol, Inglês, Química,
Física e Matemática, respectivamente.

\begin{longtable}[]{@{}cccccc@{}}
\caption{Matriz de correlações
hipotética}\label{tbl-correlacao}\tabularnewline
\toprule\noalign{}
& \(Y_1\) & \(Y_2\) & \(Y_3\) & \(Y_4\) & \(Y_5\) \\
\midrule\noalign{}
\endfirsthead
\toprule\noalign{}
& \(Y_1\) & \(Y_2\) & \(Y_3\) & \(Y_4\) & \(Y_5\) \\
\midrule\noalign{}
\endhead
\(Y_1\) & 1 & 0,90 & 0,05 & 0,05 & 0,10 \\
\(Y_2\) & 0,90 & 1 & 0,05 & 0,10 & 0,05 \\
\(Y_3\) & 0,05 & 0,05 & 1 & 0,95 & 0,85 \\
\(Y_4\) & 0,05 & 0,10 & 0,95 & 1 & 0,90 \\
\(Y_5\) & 0,10 & 0,05 & 0,85 & 0,90 & 1 \\
\bottomrule\noalign{}
\end{longtable}
\end{frame}

\begin{frame}{Análise Fatorial}
\phantomsection\label{anuxe1lise-fatorial-2}
Com base nas correlações, poderíamos imaginar dois fatores subjacentes,
e não observáveis:

\begin{itemize}
\tightlist
\item
  \textbf{Fator 1:} Conhecimento em línguas - associado aos desempenhos
  em Espanhol e Inglês, fortemente correlacionados;
\end{itemize}

\pause

\begin{itemize}
\tightlist
\item
  \textbf{Fator 2:} Conhecimento em ciências exatas - associado aos
  desempenhos em Química, Física e Matemática.
\end{itemize}

\pause

É importante compreender que ``conhecimento'' é algo que não se mede
diretamente, mas que pode se manifestar através de variáveis
observáveis, como as notas obtidas nos exames.
\end{frame}

\begin{frame}{Análise Fatorial}
\phantomsection\label{anuxe1lise-fatorial-3}
Pode ser considerada uma extensão da Análise de Componentes Principais

\pause

\begin{itemize}
\tightlist
\item
  Ambas são tentativas de aproximar \(\boldsymbol{S}\).
\end{itemize}

\pause

\begin{itemize}
\tightlist
\item
  A aproximação baseada em Análise Fatorial é mais elaborada
\end{itemize}

\pause

Questão principal:

\begin{itemize}
\tightlist
\item
  Dados são consistentes com a estrutura prescrita?
\end{itemize}
\end{frame}

\begin{frame}{Tipos de Análise Fatorial}
\phantomsection\label{tipos-de-anuxe1lise-fatorial}
\begin{block}{Análise Fatorial Exploratória}
\phantomsection\label{anuxe1lise-fatorial-exploratuxf3ria}
\medskip

\begin{itemize}
\item
  Busca encontrar os fatores subjacentes às variáveis originais
  amostradas
\item
  Em geral, efetuada quando não se tem noção clara da quantidade de
  fatores do modelo e nem do que representam
\end{itemize}
\end{block}

\begin{block}{Análise Fatorial Confirmatória}
\phantomsection\label{anuxe1lise-fatorial-confirmatuxf3ria}
\begin{itemize}
\tightlist
\item
  Tem-se em mãos um modelo fatorial pré-especificado (modelo hipotético)
  e deseja-se verificar se é aplicável ou consistente com os dados
  amostrais de que dispõe
\end{itemize}
\end{block}
\end{frame}

\begin{frame}{Planejamento de uma análise fatorial}
\phantomsection\label{planejamento-de-uma-anuxe1lise-fatorial}
\begin{itemize}
\tightlist
\item
  \textbf{Verificar a adequabilidade da base de dados:} tamanho da
  amostra (desejável a razão de 5 indivíduos para cada variável),
  correlações significativas, KMO adequado.
\end{itemize}

\pause

\begin{itemize}
\tightlist
\item
  \textbf{Determinar o método de extração e o número de fatores a serem
  extraídos:} métodos de extração (componentes principais, fatores
  principais, máxima verossimilhança), número de fatores (critério de
  kaiser, screeplot, variância acumulada).
\end{itemize}

\pause

\begin{itemize}
\tightlist
\item
  \textbf{Decidir o tipo de rotação:} Varimax, Quartimax,
  Orthomax\ldots{}
\end{itemize}
\end{frame}

\begin{frame}{Planejamento de uma análise fatorial}
\phantomsection\label{planejamento-de-uma-anuxe1lise-fatorial-1}
\begin{block}{GIGO: Garbage In, Garbage Out!}
\phantomsection\label{gigo-garbage-in-garbage-out}
\begin{tcolorbox}[enhanced jigsaw, bottomtitle=1mm, opacitybacktitle=0.6, toprule=.15mm, colframe=quarto-callout-warning-color-frame, titlerule=0mm, bottomrule=.15mm, title=\textcolor{quarto-callout-warning-color}{\faExclamationTriangle}\hspace{0.5em}{Cuidado!}, colbacktitle=quarto-callout-warning-color!10!white, breakable, colback=white, left=2mm, coltitle=black, opacityback=0, rightrule=.15mm, toptitle=1mm, arc=.35mm, leftrule=.75mm]

\begin{itemize}
\tightlist
\item
  A análise fatorial sempre produzirá fatores. Desta forma, ela é sempre
  um candidato para o fenômeno \textbf{``lixo dentro, lixo fora''}.
\end{itemize}

\medskip

\begin{itemize}
\tightlist
\item
  Se o pesquisador incluir indiscriminadamente um grande número de
  variáveis e esperar que a análise fatorial ``faça revelações'', a
  possibilidade de resultados pobres será alta.
\end{itemize}

\medskip

\begin{itemize}
\tightlist
\item
  A qualidade e o significado dos fatores determinados refletem
  \textbf{estruturas conceituais} das variáveis incluídas na análise
  (base conceitual).
\end{itemize}

\end{tcolorbox}
\end{block}
\end{frame}

\begin{frame}{Planejamento de uma análise fatorial}
\phantomsection\label{planejamento-de-uma-anuxe1lise-fatorial-2}
\begin{block}{Questões conceituais\ldots{}}
\phantomsection\label{questuxf5es-conceituais}
\bigskip

\begin{itemize}
\tightlist
\item
  É responsabilidade do pesquisador garantir que os padrões observados
  sejam conceitualmente válidos. A AF é capaz somente de determinar a
  adequação das correlações entre variáveis.
\end{itemize}

\medskip

\begin{itemize}
\tightlist
\item
  O pesquisador deve garantir que a amostra esteja adequada à estrutura
  fatorial inerente.
\end{itemize}

\medskip

\begin{itemize}
\tightlist
\item
  Por exemplo, ao avaliar os fatores de risco de uma acidente vascular,
  se nenhuma variável relacionada aos fatores genéticos for incluída, a
  AF não será capaz de identificar esta dimensão.
\end{itemize}
\end{block}
\end{frame}

\begin{frame}{Planejamento de uma análise fatorial}
\phantomsection\label{planejamento-de-uma-anuxe1lise-fatorial-3}
\begin{block}{Questões conceituais\ldots{}}
\phantomsection\label{questuxf5es-conceituais-1}
\begin{itemize}
\tightlist
\item
  A qualidade e o significado dos fatores obtidos refletem as bases
  conceituais das variáveis incluídas na análise.
\end{itemize}

\medskip

\begin{itemize}
\tightlist
\item
  O pesquisador não deve incluir indiscriminadamente inúmeras variáveis
  e esperar que a AF ``arrume as coisas''.
\end{itemize}

\medskip

\begin{itemize}
\tightlist
\item
  No exemplo da avaliação dos fatores de risco de uma acidente vascular,
  incluir a variável, por exemplo, ``cor dos olhos'' interfere nos
  resultados da AF e não traz nenhuma informação relevante para o objeto
  de estudo.
\end{itemize}
\end{block}
\end{frame}

\begin{frame}{Planejamento de uma análise fatorial}
\phantomsection\label{planejamento-de-uma-anuxe1lise-fatorial-4}
\begin{block}{Questões estatísticas\ldots{}}
\phantomsection\label{questuxf5es-estatuxedsticas}
\begin{itemize}
\tightlist
\item
  Espera-se que haja correlações significativas, pois o objetivo é
  identificar conjuntos de variáveis inter-relacionadas.
\end{itemize}

\vspace{0.3cm}

\begin{itemize}
\tightlist
\item
  Se todas as correlações são pequenas ou iguais, o pesquisador deve
  questionar o uso da análise fatorial.
\end{itemize}

\vspace{0.3cm}

\begin{itemize}
\tightlist
\item
  Inspeção visual da matriz de correlação -- coeficientes pelo menos em
  torno de 0,30.
\end{itemize}
\end{block}
\end{frame}

\begin{frame}{Planejamento de uma análise fatorial}
\phantomsection\label{planejamento-de-uma-anuxe1lise-fatorial-5}
\begin{itemize}
\tightlist
\item
  \textbf{Teste de esfericidade de Bartlett (BTS):} O teste de
  esfericidade de Bartlett testa a hipótese de que as variáveis não
  sejam correlacionadas na população.
\end{itemize}

\[H_0: \boldsymbol{P} = \boldsymbol{I}\]

A estatística do teste é dada por

\[\chi^2 = -\left[ (n-1)-\dfrac{2p + 5}{6}\right] \ln|\boldsymbol{R}|\]

que tem uma distribuição qui-quadrado com \(\nu\) graus de liberdade,
sendo \(\nu = \dfrac{p(p-1)}{2}\), \(n\) igual ao tamanho da amostra,
\(p\) igual ao número de variáveis e \(|\boldsymbol{R}|\), o
determinante da matriz de correlação amostral.
\end{frame}

\begin{frame}{Planejamento de uma análise fatorial}
\phantomsection\label{planejamento-de-uma-anuxe1lise-fatorial-6}
\begin{itemize}
\tightlist
\item
  \textbf{Medida de adequação da amostra:} comumente utilizado o KMO
  (Kaiser-Meyer-Olkin). Rencher(2002) sugere que para um modelo de
  Análise Fatorial possa ser ajustado adequadamente aos dados é
  necessário que \(\boldsymbol{R}^{-1}\) seja próxima da matriz
  diagonal. O coeficiente KMO baseia-se nesse princípio
\end{itemize}

\[KMO = \dfrac{\displaystyle{\sum_{i=1}^p}\displaystyle{\sum_{j=1}^p} r^2_{ij}}{\displaystyle{\sum_{i=1}^p}\displaystyle{\sum_{j=1}^p} r_{ij}^2 + \displaystyle{\sum_{i=1}^p}\displaystyle{\sum_{j=1}^p}v_{ij}^2},\]
\end{frame}

\begin{frame}{Planejamento de uma análise fatorial}
\phantomsection\label{planejamento-de-uma-anuxe1lise-fatorial-7}
em que \(r_{ij}\) e \(v_{ij}\) são, respectivamente, os elementos na
posição \((i,j)\) da matriz de correlações amostral, \(\boldsymbol{R}\),
e da matriz
\(\boldsymbol{V} = \boldsymbol{U} \boldsymbol{R}^{-1} \boldsymbol{U}\),
na qual
\(\boldsymbol{U} = \left[{\rm diag} (\boldsymbol{R}^{-1})^{\frac{1}{2}}\right]^{-1}\).
Note que \({\rm diag}(\boldsymbol{R}^{-1})\) é a matriz cuja diagonal
coincide com a diagonal de \(\boldsymbol{R}^{-1}\) mas com os demais
elementos nulos e \({\rm diag}(\boldsymbol{R}^{-1})^{\frac{1}{2}}\) é a
matriz diagonal cujo \(i\)-ésimo elemento diagonal é a raiz quadrada do
\(i\)-ésimo elemento diagonal de \({\rm diag}(\boldsymbol{R}^{-1})\). A
matriz \(\boldsymbol{V}\) é usualmente designada por matriz de
correlação anti-imagem.
\end{frame}

\begin{frame}{Planejamento de uma análise fatorial}
\phantomsection\label{planejamento-de-uma-anuxe1lise-fatorial-8}
\begin{itemize}
\tightlist
\item
  \textbf{Medida de adequação da amostra:} Para cada variável: MSA
  (Measure of sampling adequacy). Essa medida é bastante similar ao
  \(KMO\). Novamente, desejamos verificar a possibilidade de existência
  de uma estrutura fatorial nos dados. A \(MSA\) deve ser calculada
  separadamente para cada variável, através de
\end{itemize}

\[MSA_i = \dfrac{\displaystyle{\sum_{j=1}^p} r^2_{ij}}{\displaystyle{\sum_{j=1}^p} r_{ij}^2 + \displaystyle{\sum_{j=1}^p}v_{ij}^2}.\]
\end{frame}

\begin{frame}{Planejamento de uma análise fatorial}
\phantomsection\label{planejamento-de-uma-anuxe1lise-fatorial-9}
\begin{itemize}
\tightlist
\item
  O objetivo é verificar se uma dada variável pode ser explicada pelas
  demais, o que é esperado em um modelo fatorial.
\end{itemize}

\pause

\begin{itemize}
\tightlist
\item
  Valores baixos de \(MSA_i\) são indícios de que a respectiva variável
  pode ser retirada da análise, sem maiores prejuízos.
\end{itemize}

\pause

\begin{itemize}
\tightlist
\item
  Note que esta medida é bastante similar ao \(KMO\), apresentando uma
  lógica semelhante.
\end{itemize}
\end{frame}

\begin{frame}{Planejamento de uma análise fatorial}
\phantomsection\label{planejamento-de-uma-anuxe1lise-fatorial-10}
Ambos variam de 0 a 1, sendo:

\begin{longtable}[]{@{}cl@{}}
\caption{Tabela de interpretação do índice
KMO/MSA}\label{tbl-kmo}\tabularnewline
\toprule\noalign{}
KMO/MSA & Interpretação \\
\midrule\noalign{}
\endfirsthead
\toprule\noalign{}
KMO/MSA & Interpretação \\
\midrule\noalign{}
\endhead
\(0,80 - 1,00\) & Excelente \\
\(0,70 - 0,80\) & Ótimo \\
\(0,60 - 0,70\) & Bom \\
\(0,50 - 0,60\) & Regular \\
\(0,00 - 0,50\) & Insuficiente \\
\bottomrule\noalign{}
\end{longtable}
\end{frame}

\begin{frame}{O modelo fatorial ortogonal via matriz de correlações}
\phantomsection\label{o-modelo-fatorial-ortogonal-via-matriz-de-correlauxe7uxf5es}
\begin{itemize}
\tightlist
\item
  Seja o vetor aleatório \(\mathbf{x} = [X_1, X_2, \cdots, X_p]^t\), com
  vetor de médias \(\boldsymbol{\mu}\), matriz de covariâncias
  \(\boldsymbol{\Sigma}\) e matriz de correlações \(\boldsymbol{P}\).
\end{itemize}

\pause

\begin{itemize}
\tightlist
\item
  Sejam as variáveis originais padronizadas:
  \(Z_i = \dfrac{X_i - \mu_i}{\sqrt{\sigma_{ii}}}\)
\end{itemize}

\pause

\begin{itemize}
\tightlist
\item
  \(\boldsymbol{P}\) é a matriz de covariâncias do velor aleatório
  \(\boldsymbol{z}\), cujas componentes são as variáveis padronizadas.
\end{itemize}
\end{frame}

\begin{frame}{O modelo fatorial ortogonal via matriz de correlações}
\phantomsection\label{o-modelo-fatorial-ortogonal-via-matriz-de-correlauxe7uxf5es-1}
\begin{itemize}
\tightlist
\item
  Modelo fatorial ortogonal

  \begin{itemize}
  \tightlist
  \item
    Construído via a matriz de correlação populacional
  \item
    Relaciona linearmente as variáveis padronizadas e os \(m\) fatores
    comuns (que são desconhecidos)
  \item
    Fatores são variáveis independentes
  \end{itemize}
\end{itemize}
\end{frame}

\begin{frame}{O modelo fatorial ortogonal via matriz de correlações}
\phantomsection\label{o-modelo-fatorial-ortogonal-via-matriz-de-correlauxe7uxf5es-2}
Na análise fatorial, vamos expressar cada uma das \(p\) variáveis
padronizadas como combinação linear de \(m < p\) fatores comuns
\(F_1, F_2, \cdots, F_m\) a menos de um termo correspondente ao erro
\(\epsilon_1, \epsilon_2, \cdots, \epsilon_p\):

\[
\begin{aligned}
Z_1 &= l_{11}F_1 + l_{12}F_2 + \cdots + l_{1m}F_m + \epsilon_1 \\
Z_2 &= l_{21}F_1 + l_{22}F_2 + \cdots + l_{2m}F_m + \epsilon_2 \\ && \hspace{-25cm} \vdots \\
Z_p &= l_{p1}F_1 + l_{p2}F_2 + \cdots + l_{pm}F_m + \epsilon_p 
\end{aligned}
\]
\end{frame}

\begin{frame}{O modelo fatorial ortogonal}
\phantomsection\label{o-modelo-fatorial-ortogonal}
\begin{itemize}
\tightlist
\item
  \(F_1, F_2, \cdots, F_m\) são os fatores comuns \textbf{não
  observáveis}, ou (\emph{latentes});
\end{itemize}

\vspace{0.5cm}

\begin{itemize}
\tightlist
\item
  \(\epsilon_1, \epsilon_2, \cdots,\epsilon_p\) são os fatores
  específicos associados a cada variável e,
\end{itemize}

\vspace{0.5cm}

\begin{itemize}
\tightlist
\item
  \(l_{ij}\), \(i = 1, 2, \cdots,p\), \(j = 1, 2, \cdots, m\) são
  chamados de cargas fatoriais, ou (\emph{loadings}).
\end{itemize}
\end{frame}

\begin{frame}{O modelo fatorial ortogonal}
\phantomsection\label{o-modelo-fatorial-ortogonal-1}
\begin{block}{Notação Matricial}
\phantomsection\label{notauxe7uxe3o-matricial}
\[\boxed{\boldsymbol{D}^{-1}(\mathbf{x} - \boldsymbol{\mu}) = \boldsymbol{L} \boldsymbol{F} + \boldsymbol{\epsilon}}\]

\[\boldsymbol{F} = \left[\begin{array}{c} F_1 \\ F_2 \\ \vdots \\ F_m \end{array} \right] \boldsymbol{\epsilon} = 
\left[\begin{array}{c} \epsilon_1 \\ \epsilon_2 \\ \vdots \\ \epsilon_p \end{array}\right]
\boldsymbol{L} = \left[ \begin{array}{cccc} l_{11} & l_{12} & \cdots & l_{1m} 
\\ l_{21} & l_{22} & \cdots & l_{2m} 
\\ \vdots & \vdots  & \ddots & \vdots 
\\ l_{p1} & l_{p2} & \cdots & l_{pm}
 \end{array} \right] \,\,\, \boldsymbol{D} = \left[ \begin{array}{cccc} \sigma_{1} & 0 & \cdots & 0 
\\ 0 & \sigma_{2} & \cdots & 0 
\\ \vdots & \vdots  & \ddots & \vdots 
\\ 0 & 0 & \cdots & \sigma_{p}
 \end{array} \right]\]
\end{block}
\end{frame}

\begin{frame}{O modelo fatorial ortogonal}
\phantomsection\label{o-modelo-fatorial-ortogonal-2}
\begin{block}{Suposições do modelo}
\phantomsection\label{suposiuxe7uxf5es-do-modelo}
Para construção do modelo fatorial ortogonal, algumas suposições se
fazem necessárias:

\vspace{0.5cm}

\begin{enumerate}
[a)]
\tightlist
\item
  Os fatores têm média igual a zero, variâncias iguais a 1 e não são
  correlacionados:
\end{enumerate}

\[\boxed{E({\boldsymbol{F}}) = \boldsymbol{0}\,\, \text{e} \,\, \rm{Var}({\boldsymbol{F}}) = E({\boldsymbol{F}}{\boldsymbol{F}}^t)  = \boldsymbol{I} }\]
\end{block}
\end{frame}

\begin{frame}{O modelo fatorial ortogonal}
\phantomsection\label{o-modelo-fatorial-ortogonal-3}
\begin{block}{Suposições do modelo}
\phantomsection\label{suposiuxe7uxf5es-do-modelo-1}
\begin{enumerate}
[a)]
\setcounter{enumi}{1}
\tightlist
\item
  Os erros têm média zero, são não correlacionados e não necessariamente
  tem a mesma variância:
\end{enumerate}

\[\boxed{E({\boldsymbol{\epsilon}}) = \boldsymbol{0} e  \rm{Var}({\boldsymbol{\epsilon}}) = \boldsymbol{\Psi} = \rm{diag}(\psi_1, \psi_2, \cdots, \psi_p)
= E({\boldsymbol{\epsilon}}{\boldsymbol{\epsilon}}^t)
}\]

\begin{enumerate}
[a)]
\setcounter{enumi}{2}
\tightlist
\item
  \({\boldsymbol{F}}\) e \({\boldsymbol{\epsilon}}\) são independentes,
  isto é, fontes de variação distintas:
\end{enumerate}

\[\boxed{\rm{Cov}({\boldsymbol{F}},{\boldsymbol{\epsilon}}) = E({\boldsymbol{\epsilon}}{\boldsymbol{F}}^t) = \boldsymbol{0}}\]
\end{block}
\end{frame}

\begin{frame}{Propriedades do modelo fatorial}
\phantomsection\label{propriedades-do-modelo-fatorial}
Como consequência da definição do modelo e das suposições apresentadas,
decorrem as seguintes propriedades:

\[\boldsymbol{P} = \boldsymbol{L}\boldsymbol{L}^t + \boldsymbol{\Psi}\]

\pause

\begin{itemize}
\tightlist
\item
  Esta decomposição implica que as variâncias das variáveis observadas
  podem ser escritas como:
\end{itemize}

\[\text{Var}(Z_i) = \underbrace{l_{i1}^2 + l_{i2}^2 + \cdots + l_{im}^2}_{h_i^2} + \psi_i = h_i^2 + \psi_i, \hspace{0.5cm} i = 1, \cdots, p\]
\end{frame}

\begin{frame}{Propriedades do modelo fatorial}
\phantomsection\label{propriedades-do-modelo-fatorial-1}
\begin{itemize}
\tightlist
\item
  \(h_i^2\) é a \textbf{comunalidade:} variância comum de
  \(\boldsymbol{z}\) - expressa o quanto da variabiliade de \(Z_i\) é
  explicada pelo modelo
\end{itemize}

\vspace{0.2cm}

\begin{itemize}
\tightlist
\item
  \(\psi_i\) é a \textbf{especificidade:} variância específica de cada
  \(Z_i\) - expressa o quanto da variabilidade de \(Z_i\) não é
  explicada pelo modelo.
\end{itemize}

\pause

Além disso, as covariâncias das variáveis observadas podem ser escritas
como:

\vspace{-0.3cm}

\[\text{Cov}(Z_i, Z_k) = {l_{i1}l_{k1} + l_{i2}l_{k2} + \cdots + l_{im}l_{km}}, \,\,\, i,k = 1, \cdots, p \,\,(i\neq k)\]

\pause

e,

\[\text{Cov}(\boldsymbol{z}, \boldsymbol{F}) = \text{Cor}(\boldsymbol{z}, \boldsymbol{F}) = \boldsymbol{L}\]
\end{frame}

\begin{frame}{Métodos de Estimação}
\phantomsection\label{muxe9todos-de-estimauxe7uxe3o}
\begin{itemize}
\tightlist
\item
  \textbf{Método das componentes principais:} baseia-se na análise de
  componentes principais. Em geral, utilizado como um análise
  exploratória dos dados, em termos dos fatores subjacentes
\end{itemize}

\pause

\begin{itemize}
\tightlist
\item
  \textbf{Método dos fatores principais:} refina a solução apresentada
  pelo método das componentes principais
\end{itemize}

\pause

\begin{itemize}
\tightlist
\item
  \textbf{Método da máxima verossimilhança:} supõe que os dados seguem
  uma distribuição normal multivariada
\end{itemize}
\end{frame}

\begin{frame}{Métodos de estimação: Componentes Principais}
\phantomsection\label{muxe9todos-de-estimauxe7uxe3o-componentes-principais}
Considerando \(m < p\) o número de fatores comuns, a matriz de cargas
fatoriais, estimada pelo método dos componentes principais, fica dada
por:

\vspace{-0.3cm}

\[
\begin{aligned}
\hat{\boldsymbol{L}} &= \left[\begin{array}{cccc} \sqrt{\hat{\lambda}}_1 \hat{\boldsymbol{e}}_1 & \sqrt{\hat{\lambda}}_2 \hat{\boldsymbol{e}}_2 & \cdots & \sqrt{\hat{\lambda}}_m \hat{\boldsymbol{e}}_m \end{array}\right] 
\end{aligned}
\]

\pause

E a matriz de especificidades é dada por:

\[\hat{\boldsymbol{\Psi}} = \rm{diag}(\boldsymbol{R} - \hat{\boldsymbol{L}} \hat{\boldsymbol{L}}^t)\]

\pause

Aproximação de \(\boldsymbol{R}\)

\[\boldsymbol{R} \approx  \hat{\boldsymbol{L}} \hat{\boldsymbol{L}}^t +  \hat{\boldsymbol{\Psi}}\]
\end{frame}

\begin{frame}{Métodos de estimação: Componentes Principais}
\phantomsection\label{muxe9todos-de-estimauxe7uxe3o-componentes-principais-1}
Matriz residual:

\[\boldsymbol{MRes} = \boldsymbol{R} - (\hat{\boldsymbol{L}} \hat{\boldsymbol{L}}^t + \hat{\boldsymbol{\Psi}})\]

\pause

\begin{itemize}
\tightlist
\item
  Pode servir como critério de avaliação do modelo

  \begin{itemize}
  \tightlist
  \item
    Seus valores deveriam ser próximos de zero
  \item
    Matriz é nula somente quando o valor de \(m\) é igual a \(p\)
  \end{itemize}
\end{itemize}

\pause

\begin{itemize}
\tightlist
\item
  Os elementos da diagonal da matriz \(\boldsymbol{R}\) são reproduzidos
  exatamente pelo modelo

  \begin{itemize}
  \tightlist
  \item
    O mesmo não ocorre para os outros elementos da matriz
    \(\boldsymbol{R}\) (covariâncias das variáveis \(Z_i\) e \(Z_j\))
  \end{itemize}
\end{itemize}
\end{frame}

\begin{frame}{Métodos de estimação: Componentes Principais}
\phantomsection\label{muxe9todos-de-estimauxe7uxe3o-componentes-principais-2}
Proporção da variabilidade total explicada pelo fator

\[\text{Proporção explicada}_{F_j} = \dfrac{\sum \limits_{i=1}^p l_{ij}^2}{p}\]

\pause

\begin{itemize}
\tightlist
\item
  Representa o quanto cada fator consegue captar da variabilidade
  original das variáveis \(Z_i\).
\end{itemize}
\end{frame}

\begin{frame}{Métodos de estimação: Fatores Principais}
\phantomsection\label{muxe9todos-de-estimauxe7uxe3o-fatores-principais}
\begin{itemize}
\tightlist
\item
  Também chamado Método de Componentes Principais Iterativo
\end{itemize}

\pause

\begin{itemize}
\tightlist
\item
  \textbf{Ideia básica:} Refinar as estimativas de
  \(\hat{\boldsymbol{L}}\) e \(\hat{\boldsymbol{\Psi}}\), encontradas a
  partir do método das componentes principais.
\end{itemize}

\pause

\begin{itemize}
\tightlist
\item
  \textbf{Procedimento:}

  \begin{itemize}
  \tightlist
  \item
    Estimativas iniciais pelo método das componentes principais
  \item
    Troca dos elementos da diagonal de \(\boldsymbol{R}\) pelas
    comunalidades estimadas
  \item
    Novas estimações a partir da matriz \(\boldsymbol{R}^*\)
  \item
    Substituição dos elementos da diagonal principal pelas comunalidades
    estimadas
  \item
    Procedimento é repetido até que as diferenças entres as
    comunalidades estimadas sejam desprezíveis
  \end{itemize}
\end{itemize}
\end{frame}

\begin{frame}{Métodos de estimação: Máxima Verossimilhança}
\phantomsection\label{muxe9todos-de-estimauxe7uxe3o-muxe1xima-verossimilhanuxe7a}
Só pode ser utilizado quando a forma da distribuição de probabilidades é
conhecida

\bigskip

\begin{itemize}
\tightlist
\item
  \textbf{Suposição:}

  \begin{itemize}
  \tightlist
  \item
    Vetor aleatório \(\boldsymbol{z}\) tem distribuição normal
    \(p\)-variada
  \end{itemize}
\end{itemize}

\pause

\begin{itemize}
\tightlist
\item
  \textbf{Consequência:}

  \begin{itemize}
  \tightlist
  \item
    Vetor das variáveis padronizadas é normal \(p\)-variado
  \item
    Fatores tem distribuição normal multivariada com vetor de médias
    zero e matriz de covariâncias \(\boldsymbol{I}_m\)
  \item
    Erros tem distribuição normal \(p\)-variada com vetor de médias zero
    e matriz de covariâncias \(\boldsymbol{\Psi}\).
  \end{itemize}
\end{itemize}
\end{frame}

\begin{frame}{Métodos de estimação: Máxima Verossimilhança}
\phantomsection\label{muxe9todos-de-estimauxe7uxe3o-muxe1xima-verossimilhanuxe7a-1}
A função de verossimilhança considerando uma amostra aleatória de
tamanho \(n\) observada no vetor aleatório \(\boldsymbol{z}\) é dada
por:

\[
\boldsymbol{L}(\boldsymbol{0},\boldsymbol{P}) =  \displaystyle \frac{1}{(2\pi)^{\frac{np}{2}} |(\boldsymbol{L} \boldsymbol{L}^t + \boldsymbol{\Psi})|^{\frac{n}{2}}} \exp \left\{-\frac{1}{2} \displaystyle \sum_{j=1}^n \boldsymbol{z}_j^t (\boldsymbol{L}\boldsymbol{L}^t + \boldsymbol{\Psi})^{-1} \boldsymbol{z}_j  \right\} 
\]

\pause

\begin{itemize}
\tightlist
\item
  A função de verossimilhança depende da matrizes \(\boldsymbol{L}\) e
  \(\boldsymbol{\Psi}\), através da matriz de correlação
  \(\boldsymbol{P}\).
\end{itemize}
\end{frame}

\begin{frame}{Métodos de estimação: Máxima Verossimilhança}
\phantomsection\label{muxe9todos-de-estimauxe7uxe3o-muxe1xima-verossimilhanuxe7a-2}
\begin{itemize}
\tightlist
\item
  Os estimadores de máxima verossimilhança de \(\boldsymbol{L}\) e
  \(\boldsymbol{\Psi}\) são as matrizes \(\hat{\boldsymbol{L}}\) e
  \(\hat{\boldsymbol{\Psi}}\) que \textbf{maximizam} a função de
  verossimilhança \(\boldsymbol{L}(\boldsymbol{0},\boldsymbol{P})\).
\end{itemize}

\pause

\begin{itemize}
\tightlist
\item
  Maximização é feita por métodos numéricos.
\end{itemize}

\pause

\begin{itemize}
\tightlist
\item
  Método mais sofisticado que os métodos de componentes e fatores
  principais: produz estimativas mais precisas
\end{itemize}
\end{frame}

\begin{frame}{Métodos de estimação: Máxima Verossimilhança}
\phantomsection\label{muxe9todos-de-estimauxe7uxe3o-muxe1xima-verossimilhanuxe7a-3}
\begin{tcolorbox}[enhanced jigsaw, bottomtitle=1mm, opacitybacktitle=0.6, toprule=.15mm, colframe=quarto-callout-warning-color-frame, titlerule=0mm, bottomrule=.15mm, title=\textcolor{quarto-callout-warning-color}{\faExclamationTriangle}\hspace{0.5em}{Cuidado!}, colbacktitle=quarto-callout-warning-color!10!white, breakable, colback=white, left=2mm, coltitle=black, opacityback=0, rightrule=.15mm, toptitle=1mm, arc=.35mm, leftrule=.75mm]

\begin{itemize}
\tightlist
\item
  Está fundamentado na suposição de normalidade multivariada dos vetores
  \(\boldsymbol{z}\), \(\boldsymbol{F}\) e \(\boldsymbol{\epsilon}\).

  \begin{itemize}
  \tightlist
  \item
    Apenas a normalidade do vetor \(\boldsymbol{z}\) pode ser
    investigada a priori a partir dos dados amostrais
  \item
    Fatores e erros são variáveis aleatórias não observáveis
  \end{itemize}
\end{itemize}

\end{tcolorbox}

\pause

\begin{tcolorbox}[enhanced jigsaw, bottomtitle=1mm, opacitybacktitle=0.6, toprule=.15mm, colframe=quarto-callout-tip-color-frame, titlerule=0mm, bottomrule=.15mm, title=\textcolor{quarto-callout-tip-color}{\faLightbulb}\hspace{0.5em}{Observação Importante}, colbacktitle=quarto-callout-tip-color!10!white, breakable, colback=white, left=2mm, coltitle=black, opacityback=0, rightrule=.15mm, toptitle=1mm, arc=.35mm, leftrule=.75mm]

\begin{itemize}
\tightlist
\item
  O método de estimação por Máxima Verossimilhança (ML) pode ser
  utilizado na Análise Fatorial mesmo quando os dados não seguem uma
  distribuição normal multivariada. A falta de normalidade afeta
  principalmente os testes de significância (\(\chi^2\), \(IC\)), mas
  não compromete a qualidade da extração dos fatores. ML é robusto para
  escalas Likert com 5 ou mais categorias.
\end{itemize}

\end{tcolorbox}
\end{frame}

\begin{frame}{Rotação dos fatores}
\phantomsection\label{rotauxe7uxe3o-dos-fatores}
Uma propriedade importante do modelo fatorial ortogonal é a
\textbf{falta de unicidade das cargas fatoriais:} se \(\boldsymbol{L}\)
satisfaz à relação
\(\boldsymbol{P} = \boldsymbol{L}\boldsymbol{L}^t + \boldsymbol{\Psi}\),
então \(\boldsymbol{L}^* = \boldsymbol{L} \boldsymbol{T}\), tal que
\(\boldsymbol{T} \boldsymbol{T}^t = \boldsymbol{T}^t\boldsymbol{T} = \boldsymbol{I}\)
(\(\boldsymbol{T}\) ortogonal), também satisfaz:

\[
\begin{aligned}
\boldsymbol{L}^* \boldsymbol{L}^{*t} + \boldsymbol{\Psi} &= (\boldsymbol{L} \boldsymbol{T})(\boldsymbol{L} \boldsymbol{T})^t + \boldsymbol{\Psi} = \boldsymbol{L} \boldsymbol{T}\boldsymbol{T}^t \boldsymbol{L}^t + \boldsymbol{\Psi} \\ &= \boldsymbol{L} \boldsymbol{I} \boldsymbol{L}^t + \boldsymbol{\Psi} = 
\boldsymbol{L}\boldsymbol{L}^t + \boldsymbol{\Psi } = \boldsymbol{P}
\end{aligned}
\]

\pause

A matriz de resíduos permanece a mesma (\(\hat{h}_i^2\) e
\(\hat{\Psi}_i\))

\[\boldsymbol{R} - \hat{\boldsymbol{L}} \hat{\boldsymbol{L}}^t - \hat{\boldsymbol{\Psi}} = \boldsymbol{R} - \hat{\boldsymbol{L}}^* \hat{\boldsymbol{L}}^{*t} - \hat{\boldsymbol{\Psi}}\]

\pause

\begin{itemize}
\tightlist
\item
  Do ponto de vista estatístico é irrelevante obter
  \(\hat{\boldsymbol{L}}\) ou \(\hat{\boldsymbol{L}}^*\)
\end{itemize}
\end{frame}

\begin{frame}{Rotação dos fatores}
\phantomsection\label{rotauxe7uxe3o-dos-fatores-1}
\begin{block}{Comentários}
\phantomsection\label{comentuxe1rios}
\begin{itemize}
\tightlist
\item
  Uma vez extraídos os fatores, temos a chamada \textbf{matriz de cargas
  fatoriais}. Nesta matriz, cada elemento indica o grau de
  correspondência entre a variável e o fator;
\end{itemize}

\vspace{0.3cm}

\begin{itemize}
\tightlist
\item
  Com a rotação, busca-se uma estrutura mais simples: loadings originais
  podem não ter fácil interpretação
\end{itemize}

\vspace{0.3cm}

\begin{itemize}
\tightlist
\item
  \textbf{Ideal:} encontrar um padrão de loadings tais que cada variável
  carregue-se fortemente em um único fator (com loadings moderados nos
  outros fatores)
\end{itemize}

\vspace{0.3cm}

\begin{itemize}
\tightlist
\item
  Nem sempre é possível obter esta estrutura mais simples
\end{itemize}
\end{block}
\end{frame}

\begin{frame}{Rotação dos fatores}
\phantomsection\label{rotauxe7uxe3o-dos-fatores-2}
Ao rotacionar a matriz fatorial, busca-se redistribuir a variância dos
primeiros fatores para os últimos, a fim de atingir um padrão fatorial
mais simples e teoricamente mais significativo.

\begin{longtable}[]{@{}ccc@{}}
\caption{Cargas fatoriais dos dois fatores
extraídos}\label{tbl-fatores}\tabularnewline
\toprule\noalign{}
Variáveis & Fator 01 & Fator 02 \\
\midrule\noalign{}
\endfirsthead
\toprule\noalign{}
Variáveis & Fator 01 & Fator 02 \\
\midrule\noalign{}
\endhead
\(V_1\) & \(0,50\) & \(0,82\) \\
\(V_2\) & \(0,60\) & \(0,61\) \\
\(V_3\) & \(0,91\) & \(-0,22\) \\
\(V_4\) & \(0,82\) & \(-0,30\) \\
\(V_5\) & \(0,49\) & \(-0,50\) \\
\bottomrule\noalign{}
\end{longtable}
\end{frame}

\begin{frame}{Rotação dos fatores}
\phantomsection\label{rotauxe7uxe3o-dos-fatores-3}
\[\text{ 2 agrupamentos: As cargas fatoriais, contudo, não parecem tão óbvias}\]

\begin{columns}[T]
\begin{column}{0.48\linewidth}
\includegraphics[width=0.8\linewidth,height=\textheight,keepaspectratio]{../../images/AF/fator01.jpg}
\end{column}

\begin{column}{0.48\linewidth}
\includegraphics[width=0.8\linewidth,height=\textheight,keepaspectratio]{../../images/AF/fator02.jpg}
\end{column}
\end{columns}
\end{frame}

\begin{frame}{Rotação dos fatores}
\phantomsection\label{rotauxe7uxe3o-dos-fatores-4}
\[\text{Rotação Ortogonal}\]

\begin{columns}[T]
\begin{column}{0.48\linewidth}
\begin{longtable}[]{@{}ccc@{}}
\caption{Cargas fatoriais estimadas para dois
fatores}\label{tbl-fatores2}\tabularnewline
\toprule\noalign{}
Variáveis & Fator 01 & Fator 02 \\
\midrule\noalign{}
\endfirsthead
\toprule\noalign{}
Variáveis & Fator 01 & Fator 02 \\
\midrule\noalign{}
\endhead
\(V_1\) & \(0,03\) & \(0,90\) \\
\(V_2\) & \(0,19\) & \(0,82\) \\
\(V_3\) & \(0,95\) & \(0,22\) \\
\(V_4\) & \(0,89\) & \(0,12\) \\
\(V_5\) & \(0,78\) & \(-0,12\) \\
\bottomrule\noalign{}
\end{longtable}
\end{column}

\begin{column}{0.48\linewidth}
\includegraphics[width=0.8\linewidth,height=\textheight,keepaspectratio]{../../images/AF/fator03.jpg}
\end{column}
\end{columns}
\end{frame}

\begin{frame}{Rotação dos fatores}
\phantomsection\label{rotauxe7uxe3o-dos-fatores-5}
\begin{itemize}
\tightlist
\item
  Os métodos de rotação têm como objetivo simplificar as linhas e
  colunas da matriz fatorial para facilitar a interpretação

  \begin{itemize}
  \tightlist
  \item
    Maximizar a carga de uma variável em um único fator
  \item
    Reduzir ao máximo o número de variáveis com cargas altas por fator
  \end{itemize}
\end{itemize}

\pause

\begin{itemize}
\tightlist
\item
  Alguns critérios para encontrar matriz ortogonal:

  \begin{itemize}
  \tightlist
  \item
    Varimax
  \item
    Quartimax
  \item
    Orthomax
  \end{itemize}
\end{itemize}
\end{frame}

\begin{frame}{Rotação dos fatores}
\phantomsection\label{rotauxe7uxe3o-dos-fatores-6}
\begin{itemize}
\tightlist
\item
  \textbf{Critério Varimax:} É um dos mais utilizados. Em geral, produz
  soluções mais simples
\end{itemize}

\bigskip

\pause

\begin{itemize}
\tightlist
\item
  \textbf{Critério Quartimax:} Tem tendência de gerar fatores, onde
  todas as variáveis têm loadings elevados
\end{itemize}

\bigskip

\pause

\begin{itemize}
\tightlist
\item
  \textbf{Critério Orthomax:} É uma média ponderada dos dois outros
  métodos
\end{itemize}
\end{frame}

\begin{frame}{Rotação dos fatores}
\phantomsection\label{rotauxe7uxe3o-dos-fatores-7}
\begin{block}{Comentários}
\phantomsection\label{comentuxe1rios-1}
\begin{itemize}
\tightlist
\item
  \textbf{Qualidade de ajuste:} A rotação não acrescenta nenhuma
  melhoria em relação ao ajuste original

  \begin{itemize}
  \tightlist
  \item
    Matriz residual original não é alterada pela transformação ortogonal
  \item
    Valores estimados de comunalidade e variâncias específicas
    permanecem inalterados
  \end{itemize}
\end{itemize}

\medskip

\begin{itemize}
\tightlist
\item
  \textbf{Interpretação:} Novos fatores podem ser de mais fácil
  interpretação
\end{itemize}

\medskip

\begin{itemize}
\tightlist
\item
  Quando a solução sem rotação já é de boa qualidade, não se recomenda
  rotação: Solução rotacionada pode ser pior
\end{itemize}
\end{block}
\end{frame}

\begin{frame}{Quantos fatores usar?}
\phantomsection\label{quantos-fatores-usar}
Para determinar o valor de \(m\)\ldots{}

\begin{itemize}
\tightlist
\item
  Analisar a proporção da variação total dos dados atribuída ao
  \(j\)-ésimo fator, dada por:
\end{itemize}

\[\dfrac{\hat{\lambda}_j}{p}, \,\,\,\, j = 1,2, \cdots, m, \text{ se usarmos } \boldsymbol{R}\]

\vspace{0.5cm}

\pause

\begin{itemize}
\tightlist
\item
  \textbf{Critério de Kaiser:} retenção dos fatores com autovalor
  associado superior a 1;
\end{itemize}

\vspace{0.5cm}

\pause

\begin{itemize}
\tightlist
\item
  Gráfico scree-plot
\end{itemize}

\vspace{0.5cm}

\pause

\begin{itemize}
\tightlist
\item
  Análise paralela de Horn
\end{itemize}
\end{frame}

\begin{frame}{Quantos fatores usar?}
\phantomsection\label{quantos-fatores-usar-1}
\begin{itemize}
\tightlist
\item
  Conjugar indicadores estatísticos a:

  \begin{itemize}
  \tightlist
  \item
    Interpretação prática dos resultados;
  \item
    Busca de uma solução parcimoniosa;
  \item
    Possível indicação do número de fatores segundo a teoria da área;
  \end{itemize}
\end{itemize}

\vspace{0.3cm}

\pause

\begin{itemize}
\tightlist
\item
  \textbf{Bom senso (sempre!)}.
\end{itemize}
\end{frame}

\begin{frame}{Quantos fatores usar?}
\phantomsection\label{quantos-fatores-usar-2}
\begin{block}{Comentários}
\phantomsection\label{comentuxe1rios-2}
\begin{itemize}
\tightlist
\item
  \textbf{Método de máxima verossimilhança:} Mudança de valor de \(m\)
  altera as estimativas dos loadings
\end{itemize}

\vspace{0.5cm}

\begin{itemize}
\tightlist
\item
  \textbf{Método de componentes principais:} Aumento no valor de \(m\)
  não altera os loadings para os fatores obtidos anteriormente
\end{itemize}

\vspace{0.5cm}

\begin{itemize}
\tightlist
\item
  Quando os dados provêm de distribuição normal multivariada

  \begin{itemize}
  \tightlist
  \item
    Usar método de componentes principais como análise exploratória dos
    fatores e estimação do valor provável de \(m\)
  \item
    Posteriormente, qualidade da solução inicial poderá ser melhorada
    pelo uso do método de máxima verossimilhança
  \end{itemize}
\end{itemize}
\end{block}
\end{frame}

\begin{frame}{Interpretação dos fatores}
\phantomsection\label{interpretauxe7uxe3o-dos-fatores}
\[\text{Avaliação da significância estatística baseada no tamanho da amostra}\]

\begin{longtable}[]{@{}cc@{}}
\caption{Tamanho mínimo de amostra recomendado em função da carga
fatorial}\tabularnewline
\toprule\noalign{}
Carga Fatorial & Tamanho da amostra \\
\midrule\noalign{}
\endfirsthead
\toprule\noalign{}
Carga Fatorial & Tamanho da amostra \\
\midrule\noalign{}
\endhead
0,30 & 350 \\
0,35 & 250 \\
0,40 & 200 \\
0,45 & 150 \\
0,50 & 120 \\
0,55 & 100 \\
0,60 & 85 \\
0,65 & 70 \\
0,70 & 60 \\
0,75 & 50 \\
\bottomrule\noalign{}
\end{longtable}
\end{frame}

\begin{frame}{Interpretação dos fatores}
\phantomsection\label{interpretauxe7uxe3o-dos-fatores-1}
\[\text{Interpretação de uma matriz de cargas fatoriais em 5 passos}\]

\begin{center}
\includegraphics[width=5.20833in,height=\textheight,keepaspectratio]{../../images/AF/passos.jpg}
\end{center}
\end{frame}

\begin{frame}{Interpretação dos fatores}
\phantomsection\label{interpretauxe7uxe3o-dos-fatores-2}
\[\text{Algumas situações podem ser encontradas}\]

\[\textbf{A variável tem carga cruzada}\]

\begin{longtable}[]{@{}cccc@{}}
\toprule\noalign{}
\textbf{Variável} & \textbf{Fator 01} & \textbf{Fator 02} &
\textbf{Comunalidade} \\
\midrule\noalign{}
\endhead
\(\cdots\) & \(\cdots\) & \(\cdots\) & \(\cdots\) \\
\(V_i\) & \(-0,64\) & \(-0,56\) & \(0,72\) \\
\(\cdots\) & \(\cdots\) & \(\cdots\) & \(\cdots\) \\
\bottomrule\noalign{}
\end{longtable}
\end{frame}

\begin{frame}{Interpretação dos fatores}
\phantomsection\label{interpretauxe7uxe3o-dos-fatores-3}
\[\text{Algumas situações podem ser encontradas}\]

\[\textbf{A variável não possui cargas significantes}\]

\begin{longtable}[]{@{}cccc@{}}
\toprule\noalign{}
\textbf{Variável} & \textbf{Fator 01} & \textbf{Fator 02} &
\textbf{Comunalidade} \\
\midrule\noalign{}
\endhead
\(\cdots\) & \(\cdots\) & \(\cdots\) & \(\cdots\) \\
\(V_i\) & \(-0,25\) & \(0,19\) & \(0,10\) \\
\(\cdots\) & \(\cdots\) & \(\cdots\) & \(\cdots\) \\
\bottomrule\noalign{}
\end{longtable}
\end{frame}

\begin{frame}{Interpretação dos fatores}
\phantomsection\label{interpretauxe7uxe3o-dos-fatores-4}
\[\text{Algumas situações podem ser encontradas}\]

\[\textbf{Mesmo com cargas significantes, a comunalidade é baixa}\]

\begin{longtable}[]{@{}cccc@{}}
\toprule\noalign{}
\textbf{Variável} & \textbf{Fator 01} & \textbf{Fator 02} &
\textbf{Comunalidade} \\
\midrule\noalign{}
\endhead
\(\cdots\) & \(\cdots\) & \(\cdots\) & \(\cdots\) \\
\(V_i\) & 0,00 & 0,40 & 0,16 \\
\(\cdots\) & \(\cdots\) & \(\cdots\) & \(\cdots\) \\
\bottomrule\noalign{}
\end{longtable}
\end{frame}

\begin{frame}{Interpretação dos fatores}
\phantomsection\label{interpretauxe7uxe3o-dos-fatores-5}
\[\text{Alternativas}\]

\begin{itemize}
\tightlist
\item
  \textbf{Ignorar as variáveis problemáticas} e interpretar a solução
  fatorial como ela é, observando que as variáveis em questão são
  pobremente representadas na solução fatorial;
\end{itemize}

\vspace{0.3cm}

\pause

\begin{itemize}
\tightlist
\item
  Avaliar cada variável para uma \textbf{possível eliminação} e
  reespecificar o modelo após sua retirada;
\end{itemize}

\vspace{0.3cm}

\pause

\begin{itemize}
\tightlist
\item
  Empregar um \textbf{método alternativo de rotação};
\end{itemize}

\vspace{0.3cm}

\pause

\begin{itemize}
\tightlist
\item
  \textbf{Diminuir/aumentar o número de fatores} para avaliar se a nova
  estrutura fatorial representará as variáveis problemáticas;
\end{itemize}
\end{frame}

\begin{frame}{Comentários finais sobre análise fatorial}
\phantomsection\label{comentuxe1rios-finais-sobre-anuxe1lise-fatorial}
\begin{itemize}
\tightlist
\item
  \textbf{Matriz de resíduos:} A observação da matriz de resíduos muitas
  vezes, pode indicar quando o número de fatores está superdimensionado

  \begin{itemize}
  \tightlist
  \item
    Ex.: Se \(m\) não for muito pequeno e a matriz de resíduos estiver
    próxima de zero, recomenda-se testar outras soluções para \(m\)
    menores que o valor já especificado
  \end{itemize}
\end{itemize}

\vspace{0.3cm}

\pause

\begin{itemize}
\tightlist
\item
  A análise fatorial deve ser utilizada apenas se utilizada em situações
  em que as variáveis originais são correlacionadas: Evitar soluções com
  \(m\) elevado tal que determinados fatores fiquem relacionados com uma
  única variável original
\end{itemize}

\vspace{0.3cm}

\pause

\begin{itemize}
\tightlist
\item
  Em situações em que aparecem fatores relacionados a uma única variável
  \(Z_i\) é recomendável retirar a variável \(Z_i\) e reestimar o modelo
  de análise fatorial
\end{itemize}
\end{frame}

\begin{frame}{Usos adicionais dos resultados da AFE}
\phantomsection\label{usos-adicionais-dos-resultados-da-afe}
\[\textbf{Como posso usar os resultados de uma AFE em análises subsequentes?}\]

\begin{itemize}
\tightlist
\item
  Seleção de variáveis substitutas para análise subseqüente

  \begin{itemize}
  \tightlist
  \item
    É selecionada a variável com maior carga no fator para atuar como
    \textbf{variável substituta representativa} (deve considerar o
    conhecimento que o pesquisador tem da teoria).
  \end{itemize}
\end{itemize}

\vspace{0.3cm}

\pause

\begin{itemize}
\tightlist
\item
  \textbf{Escores Fatoriais:} representação de cada indivíduo no espaço
  fatorial. A estimação dos escores fatoriais necessita de métodos
  estatíscos elaborados, sendo os mais conhecidos os métodos da
  regressão e dos mínimos quadrados ponderados.
\end{itemize}
\end{frame}

\begin{frame}{Escores fatoriais}
\phantomsection\label{escores-fatoriais}
\begin{itemize}
\tightlist
\item
  \textbf{Método dos Mínimos Quadrados Ponderados:}
\end{itemize}

\[\hat{\boldsymbol{f}} = ( \hat{\boldsymbol{L}}^t \hat{\boldsymbol{\Psi}}^{-1} \hat{\boldsymbol{L}})^{-1}  \hat{\boldsymbol{L}}^t \hat{\boldsymbol{\Psi}}^{-1} \boldsymbol{z}\]

\bigskip

\pause

\begin{itemize}
\tightlist
\item
  \textbf{Método da Regressão:} suposição de normalidade multivariada
  para \(\boldsymbol{z}\), \(\boldsymbol{F}\) e
  \(\boldsymbol{\epsilon}\)
\end{itemize}

\[\hat{\boldsymbol{f}} = \hat{\boldsymbol{L}}^t (\hat{\boldsymbol{L}} \hat{\boldsymbol{L}}^t + \hat{\boldsymbol{\Psi}})^{-1}\boldsymbol{z}\]
\end{frame}

\begin{frame}{Contexto: empresa Telco}
\phantomsection\label{contexto-empresa-telco}
\begin{itemize}
\tightlist
\item
  Operadora de telefonia móvel em cidade do interior
\item
  Pesquisa de satisfação com os clientes após um ano de operação
\item
  Questionário com \textbf{10 itens (P1--P10)} em escala de 0 a 10
\end{itemize}
\end{frame}

\begin{frame}{Contexto: empresa Telco}
\phantomsection\label{contexto-empresa-telco-1}
Os itens avaliam, por exemplo:

\begin{itemize}
\tightlist
\item
  Intensidade do sinal
\item
  Qualidade do serviço e das informações
\item
  Justiça nas tarifas
\item
  Distribuição geográfica do sinal
\item
  Plano de serviços
\item
  Desconfiança em relação à empresa
\item
  Tecnologia, atenção e atendimento
\end{itemize}
\end{frame}

\begin{frame}{Itens e variáveis}
\phantomsection\label{itens-e-variuxe1veis}
Vamos representar os 10 itens por variáveis \(X_1, \ldots, X_{10}\):

\begin{itemize}
\tightlist
\item
  \(X_1\): intensidade do sinal
\item
  \(X_2\): distribuição geográfica do sinal
\item
  \(X_3\): tecnologia adotada
\item
  \(X_4\): atendimento como um todo
\item
  \(X_5\): injustiça nas tarifas
\item
  \(X_6\): plano de serviços
\item
  \(X_7\): qualidade das informações
\item
  \(X_8\): serviço prestado ao cliente
\item
  \(X_9\): atenção dedicada ao cliente
\item
  \(X_{10}\): atendimento especificamente ao cliente
\end{itemize}
\end{frame}

\begin{frame}[fragile]{Objetivo da análise}
\phantomsection\label{objetivo-da-anuxe1lise}
\textbf{Objetivo:} \textbf{identificar dimensões latentes} da qualidade
percebida.

\pause

\begin{Shaded}
\begin{Highlighting}[]
\CommentTok{\# Instale o pacote pacman primeiro, se ainda não o tiver}
\ControlFlowTok{if}\NormalTok{ (}\SpecialCharTok{!}\FunctionTok{requireNamespace}\NormalTok{(}\StringTok{"pacman"}\NormalTok{, }\AttributeTok{quietly =} \ConstantTok{TRUE}\NormalTok{)) \{}
  \FunctionTok{install.packages}\NormalTok{(}\StringTok{"pacman"}\NormalTok{)}
\NormalTok{\}}

\CommentTok{\# Use p\_load para instalar e carregar pacotes}
\NormalTok{pacman}\SpecialCharTok{::}\FunctionTok{p\_load}\NormalTok{(tidyverse, psych, GPArotation, ggcorrplot, janitor, corrr, igraph, ggraph, ggrepel)}
\end{Highlighting}
\end{Shaded}
\end{frame}

\begin{frame}[fragile]{Lendo os dados Telco}
\phantomsection\label{lendo-os-dados-telco}
\begin{Shaded}
\begin{Highlighting}[]
\NormalTok{path }\OtherTok{\textless{}{-}} \StringTok{"https://raw.githubusercontent.com/tiagomartin/est014/refs/heads/master/dados/telco.csv"}
\NormalTok{telco }\OtherTok{\textless{}{-}} \FunctionTok{read\_csv}\NormalTok{(path, }\AttributeTok{show\_col\_types =} \ConstantTok{FALSE}\NormalTok{) }\SpecialCharTok{\%\textgreater{}\%}
  \FunctionTok{clean\_names}\NormalTok{() }

\FunctionTok{glimpse}\NormalTok{(telco)}
\end{Highlighting}
\end{Shaded}

\begin{verbatim}
Rows: 875
Columns: 10
$ x1  <dbl> 2.021634, 2.634654, 2.161853, 2.611833, 2.296529, 2.265094, 2.7789~
$ x2  <dbl> 3.417002, 4.989637, 4.369748, 5.058855, 2.652786, 4.401388, 5.0625~
$ x3  <dbl> 3.203575, 3.548522, 3.423754, 3.438870, 3.394790, 3.446864, 3.4617~
$ x4  <dbl> 4.674381, 5.167249, 5.076795, 5.145240, 4.935702, 4.960152, 5.0267~
$ x5  <dbl> 7.845196, 7.527241, 7.432190, 7.381899, 7.572431, 7.401869, 7.3539~
$ x6  <dbl> 4.474652, 5.155799, 4.955571, 5.005566, 4.857968, 4.987190, 5.0357~
$ x7  <dbl> 1.652540, 2.303569, 3.000841, 3.084402, 2.834656, 4.363538, 4.4467~
$ x8  <dbl> 3.192175, 3.652853, 4.442697, 4.541976, 4.221692, 4.371027, 4.4521~
$ x9  <dbl> 7.092796, 7.518412, 6.690318, 6.761366, 6.098145, 6.766362, 6.8473~
$ x10 <dbl> 2.882242, 3.046624, 3.350822, 3.388591, 3.197836, 3.273597, 3.3047~
\end{verbatim}

\begin{Shaded}
\begin{Highlighting}[]
\FunctionTok{dim}\NormalTok{(telco)}
\end{Highlighting}
\end{Shaded}

\begin{verbatim}
[1] 875  10
\end{verbatim}
\end{frame}

\begin{frame}[fragile]{Matriz de correlação}
\phantomsection\label{matriz-de-correlauxe7uxe3o}
\begin{Shaded}
\begin{Highlighting}[]
\NormalTok{R }\OtherTok{\textless{}{-}} \FunctionTok{cor}\NormalTok{(telco, }\AttributeTok{use =} \StringTok{"pairwise.complete.obs"}\NormalTok{)}
\NormalTok{R}
\end{Highlighting}
\end{Shaded}

\begin{verbatim}
               x1            x2            x3            x4            x5
x1   1.0000000000  0.4781493696  0.7589493789 -2.140846e-03  0.0053772224
x2   0.4781493696  1.0000000000  0.5673144741 -6.145747e-03  0.0072963197
x3   0.7589493789  0.5673144741  1.0000000000 -5.368382e-04  0.0047779539
x4  -0.0021408464 -0.0061457470 -0.0005368382  1.000000e+00 -0.5773156240
x5   0.0053772224  0.0072963197  0.0047779539 -5.773156e-01  1.0000000000
x6   0.5414366524  0.4011884586  0.6425555030  7.028403e-01 -0.4256035677
x7   0.0014174056 -0.0011068157  0.0020693251  5.251678e-01 -0.3252902914
x8   0.0075964582  0.0028010661  0.0011511906  7.496806e-03 -0.0041722500
x9   0.0001745702  0.0013298066  0.0008497782 -3.530898e-03 -0.0026125521
x10  0.0067917792  0.0001410227  0.0044244904  3.148538e-05  0.0004600118
              x6            x7           x8            x9           x10
x1   0.541436652  0.0014174056  0.007596458  0.0001745702  6.791779e-03
x2   0.401188459 -0.0011068157  0.002801066  0.0013298066  1.410227e-04
x3   0.642555503  0.0020693251  0.001151191  0.0008497782  4.424490e-03
x4   0.702840339  0.5251677868  0.007496806 -0.0035308985  3.148538e-05
x5  -0.425603568 -0.3252902914 -0.004172250 -0.0026125521  4.600118e-04
x6   1.000000000  0.3942750794  0.008252132 -0.0023037351  3.223454e-03
x7   0.394275079  1.0000000000  0.004857482 -0.0009232318  1.151786e-03
x8   0.008252132  0.0048574820  1.000000000 -0.7925700049  5.324280e-01
x9  -0.002303735 -0.0009232318 -0.792570005  1.0000000000 -4.301518e-01
x10  0.003223454  0.0011517859  0.532427996 -0.4301518161  1.000000e+00
\end{verbatim}
\end{frame}

\begin{frame}[fragile]{Visualização da matriz de correlação}
\phantomsection\label{visualizauxe7uxe3o-da-matriz-de-correlauxe7uxe3o}
\begin{Shaded}
\begin{Highlighting}[]
\FunctionTok{ggcorrplot}\NormalTok{(}
\NormalTok{R,}
\AttributeTok{type =} \StringTok{"full"}\NormalTok{,}
\AttributeTok{lab =} \ConstantTok{TRUE}\NormalTok{,}
\AttributeTok{title =} \StringTok{"Matriz de correlação – Itens Telco"}
\NormalTok{)}
\end{Highlighting}
\end{Shaded}

\pandocbounded{\includegraphics[keepaspectratio]{fatorial_files/figure-beamer/ggcorr-1.pdf}}

\begin{Shaded}
\begin{Highlighting}[]
\CommentTok{\# Blocos de correlações altas sugerem possíveis fatores latentes}
\end{Highlighting}
\end{Shaded}
\end{frame}

\begin{frame}[fragile]{Visualização da matriz de correlação}
\phantomsection\label{visualizauxe7uxe3o-da-matriz-de-correlauxe7uxe3o-1}
\begin{Shaded}
\begin{Highlighting}[]
\DocumentationTok{\#\# Variaveis aparentam se agrupar em tres grupos}
\NormalTok{tidy\_cors }\OtherTok{\textless{}{-}}\NormalTok{ telco }\SpecialCharTok{\%\textgreater{}\%} \FunctionTok{correlate}\NormalTok{() }\SpecialCharTok{\%\textgreater{}\%} \FunctionTok{stretch}\NormalTok{()}
\NormalTok{graph\_cors }\OtherTok{\textless{}{-}}\NormalTok{ tidy\_cors }\SpecialCharTok{\%\textgreater{}\%}  \FunctionTok{filter}\NormalTok{(}\FunctionTok{abs}\NormalTok{(r) }\SpecialCharTok{\textgreater{}}\NormalTok{ .}\DecValTok{3}\NormalTok{) }\SpecialCharTok{\%\textgreater{}\%} \FunctionTok{graph\_from\_data\_frame}\NormalTok{(}\AttributeTok{directed =} \ConstantTok{FALSE}\NormalTok{)}
\FunctionTok{ggraph}\NormalTok{(graph\_cors, }\AttributeTok{layout =} \StringTok{"stress"}\NormalTok{) }\SpecialCharTok{+}
  \FunctionTok{geom\_edge\_link}\NormalTok{(}\AttributeTok{alpha =} \FloatTok{0.5}\NormalTok{) }\SpecialCharTok{+}
  \FunctionTok{geom\_node\_point}\NormalTok{(}\AttributeTok{size =} \DecValTok{4}\NormalTok{) }\SpecialCharTok{+}
  \FunctionTok{geom\_node\_text}\NormalTok{(}\FunctionTok{aes}\NormalTok{(}\AttributeTok{label =}\NormalTok{ name), }\AttributeTok{repel =} \ConstantTok{TRUE}\NormalTok{, }\AttributeTok{size =} \DecValTok{5}\NormalTok{) }\SpecialCharTok{+}
  \FunctionTok{theme\_minimal}\NormalTok{()}
\end{Highlighting}
\end{Shaded}

\pandocbounded{\includegraphics[keepaspectratio]{fatorial_files/figure-beamer/unnamed-chunk-2-1.pdf}}
\end{frame}

\begin{frame}[fragile]{Medidas de adequação à AFE}
\phantomsection\label{medidas-de-adequauxe7uxe3o-uxe0-afe}
\begin{block}{KMO -- Kaiser-Meyer-Olkin}
\phantomsection\label{kmo-kaiser-meyer-olkin}
\begin{Shaded}
\begin{Highlighting}[]
\FunctionTok{KMO}\NormalTok{(R)}
\end{Highlighting}
\end{Shaded}

\begin{verbatim}
Kaiser-Meyer-Olkin factor adequacy
Call: KMO(r = R)
Overall MSA =  0.63
MSA for each item = 
  x1   x2   x3   x4   x5   x6   x7   x8   x9  x10 
0.85 0.93 0.56 0.46 0.93 0.55 0.94 0.58 0.60 0.79 
\end{verbatim}

\begin{itemize}
\tightlist
\item
  KMO geral \textgreater{} 0,60 → dados adequados para AFE
\item
  MSA de \(X_4 < 0,5\) → apresenta adequabilidade insuficiente, podendo
  ser retirada da analise
\end{itemize}
\end{block}
\end{frame}

\begin{frame}[fragile]{Medidas de adequação à AFE}
\phantomsection\label{medidas-de-adequauxe7uxe3o-uxe0-afe-1}
\begin{block}{Teste de esfericidade de Bartlett}
\phantomsection\label{teste-de-esfericidade-de-bartlett}
\begin{Shaded}
\begin{Highlighting}[]
\FunctionTok{cortest.bartlett}\NormalTok{(R, }\AttributeTok{n =} \FunctionTok{nrow}\NormalTok{(telco))}
\end{Highlighting}
\end{Shaded}

\begin{verbatim}
$chisq
[1] 5049.455

$p.value
[1] 0

$df
[1] 45
\end{verbatim}

\begin{itemize}
\tightlist
\item
  p-valor de Bartlett \textless{} 0,05 → correlações não são nulas
\end{itemize}
\end{block}
\end{frame}

\begin{frame}[fragile]{Escolha do número de fatores: Parallel Analysis}
\phantomsection\label{escolha-do-nuxfamero-de-fatores-parallel-analysis}
\begin{Shaded}
\begin{Highlighting}[]
\FunctionTok{set.seed}\NormalTok{(}\DecValTok{123}\NormalTok{)}

\FunctionTok{fa.parallel}\NormalTok{(}
\NormalTok{telco,}
\AttributeTok{fm =} \StringTok{"ml"}\NormalTok{,}
\AttributeTok{fa =} \StringTok{"fa"}\NormalTok{,}
\AttributeTok{main =} \StringTok{"Parallel Analysis – Caso Telco"}
\NormalTok{)}
\end{Highlighting}
\end{Shaded}

\pandocbounded{\includegraphics[keepaspectratio]{fatorial_files/figure-beamer/parallel-1.pdf}}

\begin{verbatim}
Parallel analysis suggests that the number of factors =  3  and the number of components =  NA 
\end{verbatim}
\end{frame}

\begin{frame}[fragile]{Escolha do número de fatores: Parallel Analysis}
\phantomsection\label{escolha-do-nuxfamero-de-fatores-parallel-analysis-1}
\begin{itemize}
\item
  Comparamos autovalores observados com autovalores obtidos ao acaso.
\item
  Mantemos os fatores com autovalores acima da linha simulada.
\end{itemize}

\begin{Shaded}
\begin{Highlighting}[]
\NormalTok{n\_fatores }\OtherTok{\textless{}{-}} \DecValTok{3}
\end{Highlighting}
\end{Shaded}
\end{frame}

\begin{frame}[fragile]{Ajuste da AFE}
\phantomsection\label{ajuste-da-afe}
\begin{Shaded}
\begin{Highlighting}[]
\NormalTok{fa\_telco }\OtherTok{\textless{}{-}} \FunctionTok{factanal}\NormalTok{(}
  \AttributeTok{x        =}\NormalTok{ telco,}
  \AttributeTok{factors  =}\NormalTok{ n\_fatores,}
  \AttributeTok{rotation =} \StringTok{"none"}\NormalTok{,}
  \AttributeTok{scores   =} \StringTok{"none"}
\NormalTok{)}

\FunctionTok{print}\NormalTok{(fa\_telco, }\AttributeTok{digits =} \DecValTok{3}\NormalTok{, }\AttributeTok{cut =} \FloatTok{0.30}\NormalTok{) }\CommentTok{\# n = 875: cargas acima de 0,30 sao significativas}
\end{Highlighting}
\end{Shaded}

\begin{verbatim}

Call:
factanal(x = telco, factors = n_fatores, scores = "none", rotation = "none")

Uniquenesses:
   x1    x2    x3    x4    x5    x6    x7    x8    x9   x10 
0.359 0.642 0.101 0.056 0.647 0.016 0.707 0.019 0.360 0.711 

Loadings:
    Factor1 Factor2 Factor3
x1   0.514           0.614 
x2   0.380           0.462 
x3   0.610           0.726 
x4   0.743          -0.626 
x5  -0.451           0.387 
x6   0.990                 
x7   0.416          -0.346 
x8           0.990         
x9          -0.800         
x10          0.537         

               Factor1 Factor2 Factor3
SS loadings      2.691   1.910   1.780
Proportion Var   0.269   0.191   0.178
Cumulative Var   0.269   0.460   0.638

Test of the hypothesis that 3 factors are sufficient.
The chi square statistic is 0.32 on 18 degrees of freedom.
The p-value is 1 
\end{verbatim}

\begin{Shaded}
\begin{Highlighting}[]
\CommentTok{\# Podemos notar que os 3 fatores explicam juntos 63,8\% da variacao total do sistema original}
\end{Highlighting}
\end{Shaded}
\end{frame}

\begin{frame}{Ajuste da AFE}
\phantomsection\label{ajuste-da-afe-1}
\textbf{Analisando as cargas significativas}

\begin{itemize}
\tightlist
\item
  \(X_1\) - Apresenta cargas cruzadas - fatores 01 e 03 - \(h^2\): 0,641
  - MSA: 0,85
\item
  \(X_2\) - Apresenta cargas cruzadas - fatores 01 e 03 - \(h^2\): 0,358
  - MSA: 0,93
\item
  \(X_3\) - Apresenta cargas cruzadas - fatores 01 e 03 - \(h^2\): 0,899
  - MSA: 0,56
\item
  \(X_4\) - Apresenta cargas cruzadas - fatores 01 e 03 - \(h^2\): 0,944
  - MSA: 0,46
\item
  \(X_5\) - Apresenta cargas cruzadas - fatores 01 e 03 - \(h^2\): 0,353
  - MSA: 0,93
\item
  \(X_6\) - fator 01 - \(h^2\): 0,984 - MSA: 0,55
\item
  \(X_7\) - Apresenta cargas cruzadas - fatores 01 e 03 - \(h^2\): 0,293
  - MSA: 0,94
\item
  \(X_8\) - fator 02 - \(h^2\): 0,981 - MSA: 0,58
\item
  \(X_9\) - fator 02 - \(h^2\): 0,640 - MSA: 0,60
\item
  \(X_{10}\) - fator 02 - \(h^2\): 0,289 - MSA: 0,79
\end{itemize}
\end{frame}

\begin{frame}[fragile]{Vamos em busca de uma solucao melhor: rotacao
Varimax()}
\phantomsection\label{vamos-em-busca-de-uma-solucao-melhor-rotacao-varimax}
\begin{Shaded}
\begin{Highlighting}[]
\FunctionTok{varimax}\NormalTok{(fa\_telco}\SpecialCharTok{$}\NormalTok{loadings)}
\end{Highlighting}
\end{Shaded}

\begin{verbatim}
$loadings

Loadings:
    Factor1 Factor2 Factor3
x1                   0.801 
x2                   0.598 
x3                   0.948 
x4   0.971                 
x5  -0.594                 
x6   0.729           0.673 
x7   0.541                 
x8           0.991         
x9          -0.800         
x10          0.538         

               Factor1 Factor2 Factor3
SS loadings      2.121   1.910   2.350
Proportion Var   0.212   0.191   0.235
Cumulative Var   0.212   0.403   0.638

$rotmat
            [,1]        [,2]        [,3]
[1,]  0.76987948 0.025103554  0.63769538
[2,] -0.01744908 0.999680498 -0.01828749
[3,] -0.63795072 0.002951969  0.77007153
\end{verbatim}

\begin{Shaded}
\begin{Highlighting}[]
\CommentTok{\# Note que a proporcao total da variacao explicada pelos 3 fatores nao muda (63,8\%). O que muda eh a proporcao da variacao total explicada por cada fator}
\end{Highlighting}
\end{Shaded}
\end{frame}

\begin{frame}{Rotação Varimax()}
\phantomsection\label{rotauxe7uxe3o-varimax}
\textbf{Analisando novamente as cargas significativas}

\begin{itemize}
\tightlist
\item
  \(X_1\) - fator 03
\item
  \(X_2\) - fator 03
\item
  \(X_3\) - fator 03
\item
  \(X_4\) - fator 01
\item
  \(X_5\) - fator 01
\item
  \(X_6\) - Apresenta cargas cruzadas - fatores 01 e 03 - \(h^2\): 0,984
  - MSA: 0,55
\item
  \(X_7\) - fator 01
\item
  \(X_8\) - fator 02
\item
  \(X_9\) - fator 02
\item
  \(X_{10}\) - fator 02
\end{itemize}
\end{frame}

\begin{frame}[fragile]{Segunda tentativa: excluir da analise a variavel
\(X_4\) (MSA inadequado)}
\phantomsection\label{segunda-tentativa-excluir-da-analise-a-variavel-x_4-msa-inadequado}
\begin{Shaded}
\begin{Highlighting}[]
\NormalTok{telco2 }\OtherTok{\textless{}{-}}\NormalTok{ telco }\SpecialCharTok{\%\textgreater{}\%} \FunctionTok{select}\NormalTok{(}\SpecialCharTok{{-}}\NormalTok{x4)}

\NormalTok{fa\_telco2 }\OtherTok{\textless{}{-}} \FunctionTok{factanal}\NormalTok{(}
  \AttributeTok{x        =}\NormalTok{ telco2,}
  \AttributeTok{factors  =}\NormalTok{ n\_fatores,}
  \AttributeTok{rotation =} \StringTok{"varimax"}\NormalTok{,}
  \AttributeTok{scores   =} \StringTok{"none"}
\NormalTok{)}

\FunctionTok{print}\NormalTok{(fa\_telco2, }\AttributeTok{digits =} \DecValTok{3}\NormalTok{, }\AttributeTok{cut =} \FloatTok{0.30}\NormalTok{) }\CommentTok{\# n = 875: cargas acima de 0,30 sao significativas}
\end{Highlighting}
\end{Shaded}

\begin{verbatim}

Call:
factanal(x = telco2, factors = n_fatores, scores = "none", rotation = "varimax")

Uniquenesses:
   x1    x2    x3    x5    x6    x7    x8    x9   x10 
0.360 0.642 0.101 0.644 0.023 0.702 0.019 0.360 0.711 

Loadings:
    Factor1 Factor2 Factor3
x1   0.800                 
x2   0.598                 
x3   0.948                 
x5                  -0.596 
x6   0.662           0.734 
x7                   0.546 
x8           0.991         
x9          -0.800         
x10          0.537         

               Factor1 Factor2 Factor3
SS loadings      2.335   1.910   1.193
Proportion Var   0.259   0.212   0.133
Cumulative Var   0.259   0.472   0.604

Test of the hypothesis that 3 factors are sufficient.
The chi square statistic is 0.23 on 12 degrees of freedom.
The p-value is 1 
\end{verbatim}

\begin{Shaded}
\begin{Highlighting}[]
\CommentTok{\# Podemos notar que os 3 fatores explicam juntos 60,4\% da variacao total do sistema original e não resolveu o problema}
\end{Highlighting}
\end{Shaded}
\end{frame}

\begin{frame}[fragile]{Terceira tentativa: voltar para a análise a
variável \(X_4\) e retirar a variável \(X_6\)}
\phantomsection\label{terceira-tentativa-voltar-para-a-anuxe1lise-a-variuxe1vel-x_4-e-retirar-a-variuxe1vel-x_6}
\begin{Shaded}
\begin{Highlighting}[]
\NormalTok{telco3 }\OtherTok{\textless{}{-}}\NormalTok{ telco }\SpecialCharTok{\%\textgreater{}\%} \FunctionTok{select}\NormalTok{(}\SpecialCharTok{{-}}\NormalTok{x6)}

\NormalTok{fa\_telco3 }\OtherTok{\textless{}{-}} \FunctionTok{factanal}\NormalTok{(}
  \AttributeTok{x        =}\NormalTok{ telco3,}
  \AttributeTok{factors  =}\NormalTok{ n\_fatores,}
  \AttributeTok{rotation =} \StringTok{"none"}\NormalTok{,}
  \AttributeTok{scores   =} \StringTok{"Bartlett"}
\NormalTok{)}

\FunctionTok{print}\NormalTok{(fa\_telco3, }\AttributeTok{digits =} \DecValTok{3}\NormalTok{, }\AttributeTok{cut =} \FloatTok{0.30}\NormalTok{) }\CommentTok{\# n = 875: cargas acima de 0,30 sao significativas}
\end{Highlighting}
\end{Shaded}

\begin{verbatim}

Call:
factanal(x = telco3, factors = n_fatores, scores = "Bartlett",     rotation = "none")

Uniquenesses:
   x1    x2    x3    x4    x5    x7    x8    x9   x10 
0.360 0.643 0.100 0.068 0.642 0.704 0.019 0.360 0.711 

Loadings:
    Factor1 Factor2 Factor3
x1           0.800         
x2           0.598         
x3           0.949         
x4                   0.965 
x5                  -0.598 
x7                   0.544 
x8   0.991                 
x9  -0.800                 
x10  0.537                 

               Factor1 Factor2 Factor3
SS loadings      1.911   1.897   1.586
Proportion Var   0.212   0.211   0.176
Cumulative Var   0.212   0.423   0.599

Test of the hypothesis that 3 factors are sufficient.
The chi square statistic is 0.24 on 12 degrees of freedom.
The p-value is 1 
\end{verbatim}

\begin{Shaded}
\begin{Highlighting}[]
\CommentTok{\# Podemos notar que os 3 fatores explicam juntos 59,9\% da variacao total do sistema original}
\end{Highlighting}
\end{Shaded}
\end{frame}

\begin{frame}{Solução final excluindo \(X_6\) da análise}
\phantomsection\label{soluuxe7uxe3o-final-excluindo-x_6-da-anuxe1lise}
\begin{itemize}
\tightlist
\item
  \textbf{Fator 01:} associado a \(X_8\), \(X_9\) e \(X_{10}\)

  \begin{itemize}
  \tightlist
  \item
    \textbf{Interpretação:} dimensão relacionada a atendimento /
    relacionamento com o cliente
  \end{itemize}
\item
  \textbf{Fator 02:} associado a \(X_1\), \(X_2\) e \(X_3\)

  \begin{itemize}
  \tightlist
  \item
    \textbf{Interpretação:} dimensão de qualidade técnica /
    infraestrutura (sinal, distribuição, tecnologia)
  \end{itemize}
\item
  \textbf{Fator 03:} associado a \(X_4\), \(X_5\) e \(X_7\)

  \begin{itemize}
  \tightlist
  \item
    \textbf{Interpretação:} dimensão de plano, informações e justiça
    tarifária
  \end{itemize}
\end{itemize}
\end{frame}

\begin{frame}[fragile]{Solução final excluindo \(X_6\) da análise}
\phantomsection\label{soluuxe7uxe3o-final-excluindo-x_6-da-anuxe1lise-1}
\begin{Shaded}
\begin{Highlighting}[]
\NormalTok{loadings\_nr }\OtherTok{\textless{}{-}}\NormalTok{ fa\_telco3}\SpecialCharTok{$}\NormalTok{loadings }\SpecialCharTok{|\textgreater{}}
\FunctionTok{unclass}\NormalTok{() }\SpecialCharTok{|\textgreater{}}
\FunctionTok{as.data.frame}\NormalTok{() }\SpecialCharTok{|\textgreater{}}
\FunctionTok{rownames\_to\_column}\NormalTok{(}\StringTok{"item"}\NormalTok{)}

\FunctionTok{ggplot}\NormalTok{(loadings\_nr, }\FunctionTok{aes}\NormalTok{(}\AttributeTok{x =}\NormalTok{ Factor1, }\AttributeTok{y =}\NormalTok{ Factor2, }\AttributeTok{label =}\NormalTok{ item)) }\SpecialCharTok{+}
\FunctionTok{geom\_hline}\NormalTok{(}\AttributeTok{yintercept =} \DecValTok{0}\NormalTok{, }\AttributeTok{linetype =} \StringTok{"dashed"}\NormalTok{, }\AttributeTok{color =} \StringTok{"grey70"}\NormalTok{) }\SpecialCharTok{+}
\FunctionTok{geom\_vline}\NormalTok{(}\AttributeTok{xintercept =} \DecValTok{0}\NormalTok{, }\AttributeTok{linetype =} \StringTok{"dashed"}\NormalTok{, }\AttributeTok{color =} \StringTok{"grey70"}\NormalTok{) }\SpecialCharTok{+}
\FunctionTok{geom\_point}\NormalTok{(}\AttributeTok{size =} \DecValTok{3}\NormalTok{) }\SpecialCharTok{+}
\FunctionTok{geom\_text\_repel}\NormalTok{(}\AttributeTok{size =} \DecValTok{5}\NormalTok{) }\SpecialCharTok{+}
\FunctionTok{labs}\NormalTok{(}
\AttributeTok{title =} \StringTok{"Mapa das cargas não rotacionadas – Fator 1 × Fator 2"}\NormalTok{,}
\AttributeTok{x =} \StringTok{"Factor 1"}\NormalTok{,}
\AttributeTok{y =} \StringTok{"Factor 2"}
\NormalTok{) }\SpecialCharTok{+}
\FunctionTok{theme\_minimal}\NormalTok{(}\AttributeTok{base\_size =} \DecValTok{14}\NormalTok{)}
\end{Highlighting}
\end{Shaded}

\pandocbounded{\includegraphics[keepaspectratio]{fatorial_files/figure-beamer/unnamed-chunk-7-1.pdf}}
\end{frame}

\begin{frame}{Solução final excluindo \(X_6\) da análise}
\phantomsection\label{soluuxe7uxe3o-final-excluindo-x_6-da-anuxe1lise-2}
\textbf{Pontos a destacar:}

\begin{itemize}
\item
  Proporção da variabilidade explicada por cada fator

  \begin{itemize}
  \tightlist
  \item
    \textbf{Fator 1:} cerca de 21\%
  \item
    \textbf{Fator 2:} cerca de 21\%
  \item
    \textbf{Fator 3:} cerca de 18\%
  \end{itemize}
\item
  Proporção da variabilidade explicada pelo modelo ≈ 0.60:
\item
  A solução de 3 fatores explica aproximadamente 60\% da variância total
  das variáveis padronizadas.
\item
  Para um estudo com itens de questionário (escala de 0 a 10), essa
  proporção é considerada muito boa.
\end{itemize}
\end{frame}




\end{document}
